\section{Counts and formulas under changes of scale}
\label{ext:scaling-specializations}

The constructions in the main text use the edge-two tetrahedron $K$.
A similarity of ratio $a/2$ produces a regular tetrahedron of edge length
$a$: it multiplies every base area and the minimum projection area by
$(a/2)^2$, while preserving the number of cylinders and the normalized
cost.  This appendix records several parameter choices and their
coordinate dictionaries.

\subsection{Secant sectors}
\label{ext:secant-specializations}

In the following table, every row uses an equal partition with $m$ intervals
and hence exactly $2m$ cylinders.  The body has edge length $a$ and is
obtained from $K$ by a similarity of ratio $a/2$; accordingly the area
is reported as $C/A_{\min}$.  Each statement holds for all sufficiently
large indicated integers or all sufficiently small indicated positive
real parameters.  Every displayed remainder is fourth order in the
displayed tilt parameter, uniformly over the permitted finer meshes.
\begin{center}
\small
\renewcommand{\arraystretch}{1.3}
\begin{tabular}{lllll}
\toprule
Cap & $a$ & $(\tau,\eta)$ & $m$ & $C/A_{\min}$\\
\midrule
Midpoint & $4$ & $(2/n,1/1920)$ & $2n^2$
 & $\frac12-\frac1{80n^2}+O(n^{-4})$\\
Midpoint & $1$ & $(\alpha,1/1000)$ & $\lceil\alpha^{-2}\rceil$
 & $\frac12-\frac{13}{6000}\alpha^2+O(\alpha^4)$\\
Midpoint & $2$ & $(1/n,1/1000)$ & $n^2$
 & $\frac12-\frac{13}{6000n^2}+O(n^{-4})$\\
Continuous & $2$ & $(c,1/1000)$ & any $m\ge2/c^2$
 & $\frac12-\frac{13}{6000}c^2+O(c^4)$\\
Squared radius & $2$ & $(2/n,1/960)$ & $n^2$
 & $\frac12-\frac1{120n^2}+O(n^{-4})$\\
Maximum & $1$ & $(1/(4n),1/960)$ & $n^2$
 & $\frac12-\frac1{7680n^2}+O(n^{-4})$\\
Cartesian, $b=3/4$ & $\sqrt2$ & $(\lambda/2,1/960)$
 & $\lceil\sqrt2/\lambda^2\rceil$
 & $\frac12-\frac{\lambda^2}{1920}+O(\lambda^4)$\\
Cartesian, $b=1$ & $1$ & $(1/(2n),1/1000)$ & $n^2$
 & $\frac12-\frac{13}{24000n^2}+O(n^{-4})$\\
\bottomrule
\end{tabular}
\end{center}
The table follows by substitution in \eqref{ext:general-cost}; in each
row the mesh is bounded by a fixed multiple of $\tau^2$.
For the continuous cap, $m=\lceil2/c^2\rceil$ gives the ceiling count,
whereas $c=1/n$ permits every integer $m\ge2n^2$, with the same
threshold and error constant.  The real-parameter midpoint row includes
the reciprocal-integer construction after dilation by two.
For the squared-radius row, the two families have equal area and each
costs $H[1/4+(1/240-1/120)n^{-2}+O(n^{-4})]$.
For the first Cartesian row, either family costs
$c_0/4-c_0\lambda^2/3840+O(\lambda^4)$ with $c_0=1/\sqrt2$.
These are genuine finite covers: the positive margin is fixed first,
then a sufficiently small tilt and the specified finite mesh are chosen.
No further limit in the number of cylinders is needed.


\subsection{Node-centred sectors at edge length \texorpdfstring{$\sqrt2$}{sqrt(2)}}
\label{ext:endpoint-scaled}

For the endpoint construction of Section~\ref{ext:endpoint-sectors},
write $t=1/k$.  Its $2(k^2+1)$ cylinders retain that exact count under
scaling.  At edge length $\sqrt2$, the similarity ratio is $1/\sqrt2$
and the minimum projection area is $H/2=1/\sqrt2$.  Multiplying
\eqref{ext:endpoint-cost} by this area gives total base area
\begin{equation}\label{ext:endpoint-scaled-cost}
 \frac{2}{\sqrt2}\left(\frac14-\frac{13}{12000}t^2+O(t^4)\right).
\end{equation}

\subsection{The Cartesian discriminant at edge length one}
\label{ext:cartesian-edge-one}

The following dictionary translates the Cartesian argument of
Appendix~\ref{ext:cartesian-direct}.  In its edge-two coordinates,
$(x,y,Hs)$ is the point of $K$, $Y_0$ is the transverse intercept,
$S$ is its radial height, and $T$ is $s$ or $1-s$ in the respective
family.  The discriminant is \eqref{ext:cartesian-discriminant}.
At edge-one scale, writing the point as $(x_1,y_1,s/\sqrt2)$ and
putting $\tau=\varepsilon/2$, the substitutions
$x=2x_1$, $y=2y_1$, $\eta=1/1000$ give the equivalent obstruction
$\varepsilon^2[4(x_1^2-y_1^2)^2+1/4000]+O(\varepsilon^3)$.
In these coordinates, put $(Y_1,\delta)=(y_1,\varepsilon x_1)$
or $(x_1,-\varepsilon y_1)$ in the two families.  The transverse
intercept is $b_1=Y_0/2=Y_1-\delta\alpha_I/2$.
For $S>0$, write $q=Y_0/S$. The edge-one slope coordinate is
$q_1=b_1/S=q/2$, whose sector has width $d_1=\Delta/2$.
The same discriminant is $Q=T^2-2\delta Y_1+\delta^2/4$, and its
bound is $\sqrt Q\ge S-|\delta|d_1/2$.
Writing $F(v)=v^2/4+4v^4$, the corresponding bounded aperture is
\[
 0\le S\le1,\qquad
 \frac l2S\le b_1\le\frac r2S,\qquad
 S\le1/2+\varepsilon^2(F(b_1)+1/4000).
\]
The outer bound $S\le1$ remains part of the definition after scaling.

\subsection{The cubic buffer at edge length four}
\label{ext:cubic-edge-four}

\begin{remark}[The edge-four normalization]\label{rem:angular-cubic-scale}
The similarity
\[
 (x,y,Ht)\longmapsto(2x,2y,H(2t-1))
\]
sends $K$ onto the regular edge-four tetrahedron
\[
 K_4=\{(x,y,Hz):-1\le z\le1,\quad
                  |x|\le1-z,\quad |y|\le1+z\}.
\]
It multiplies every area by four and leaves the cylinder count unchanged.
Writing $\varepsilon=\tau/2$, the same cover has
$N=\lceil2/\varepsilon^2\rceil$ sectors in each family. On a sector
$I=[l,r]$, write $T_I$ for its cutoff in \eqref{eq:cubic-caps}, and put
\[
 a_I=2\alpha_I=\frac{1+lr}{2},\qquad
 b_I=2\beta_I=\frac{l+r}{2},\qquad
 D_I=\max_{q\in I}\frac{q^2+q^4}{2}
     =8\max_{q\in I}d(q).
\]
The two axes are $(1,\varepsilon a_I,H\varepsilon b_I)$ and
$(-\varepsilon a_I,1,H\varepsilon b_I)$. Their radial intercept
coordinates are $1+z$ and $1-z$, with common cutoff
\[
 2T_I=1+\delta_I,\qquad
 \delta_I=\varepsilon^2D_I+100\varepsilon^3.
\]
Thus these $2\lceil2/\varepsilon^2\rceil$ triangular cylinders cover
$K_4$ for every $0<\varepsilon\le1/2$. Since
$A_{\min}(K_4)=4H$, their total area $S_\varepsilon=4C_{2\varepsilon}$
satisfies
\[
 S_\varepsilon
 =2\sqrt2-\frac{\sqrt2}{15}\varepsilon^2+O(\varepsilon^3),
\]
and is strictly below $A_{\min}(K_4)/2$ for
$0<\varepsilon\le1/8000$. More precisely,
\[
 \left|\frac{S_\varepsilon}{H}-2+\frac{\varepsilon^2}{15}
                         -400\varepsilon^3\right|
 \le128\varepsilon^4\qquad(0<\varepsilon\le1/100).
\]
The full coverage interval and the smaller sufficient saving interval
remain distinct under this similarity.
\end{remark}
